\chapter{Масштабы, струны и двадцать шесть измерений}

\lectureoverview{Что задает расстояние между уровнями струнного спектра? Ответ — натяжение, величина с размерностью квадрата массы. Для фундаментальных струн это наблюдение ведет к планковским единицам, к практической невозможности прямого возбуждения и к началу режима сильной гравитации. После намеренного отступления мы возвращаемся к основному состоянию бозонной открытой струны. Регуляризация ее энергии нулевых колебаний дает исторически \emph{простой} путь к 24 поперечным направлениям и, следовательно, к 26 измерениям пространства-времени.}

\section{Масштаб, который мы оставили в стороне}

До сих пор мы говорили о колеблющихся и вращающихся струнах, не объясняя, что означает одна единица возбуждения. Адрон можно представить как струну, и его осцилляторы могут возбуждаться, но насколько велик скачок от одного состояния к следующему? Что его определяет?

В элементарной струнной модели есть только один размерный параметр. Чтобы его найти, начнем с обычной растянутой пружины. Если ее концы раздвинуты на величину \(\ell\), закон Гука дает
\begin{equation}
  E_{\mathrm{spring}}=\frac{k\ell^2}{2}.
  \label{eq:l05-hooke}
\end{equation}
Коэффициент одна вторая был добавлен аудиторией после того, как была записана первая версия. В этой начальной оценке от таких численных множителей ничего существенного не зависит.

\begin{figure}[H]
  \centering
  \includegraphics[width=0.72\textwidth]{lecture_05_figure_01.png}
  \caption{Растянутый участок и его потенциальная энергия по закону Гука.}
  \lectureframe{00:02:28}
  \label{fig:l05-hooke-board}
\end{figure}

Энергия в \eqref{eq:l05-hooke} нерелятивистская. В описании с бесконечным импульсом, использованном в предыдущих лекциях, внутренняя нерелятивистская энергия пропорциональна не \(m\), а \(m^2\). Поэтому схематически пишем
\begin{equation}
  m^2\sim k\ell^2,
  \qquad
  m\sim \sqrt{k}\,\ell.
  \label{eq:l05-mass-length}
\end{equation}
В системе покоя струны ее масса равна ее энергии, поэтому второе соотношение принимает вид
\begin{equation}
  E_{\mathrm{rest}}\sim T_s\ell,
  \qquad
  T_s\equiv\sqrt{k}.
  \label{eq:l05-tension}
\end{equation}
Коэффициент \(T_s\) — это натяжение струны: энергия, запасенная на единицу длины.

\begin{figure}[H]
  \centering
  \includegraphics[width=0.78\textwidth]{lecture_05_figure_02.png}
  \caption{На доске переходят от \(m^2\sim k\ell^2\) к массе, пропорциональной длине.}
  \lectureframe{00:05:42}
  \label{fig:l05-mass-length-board}
\end{figure}

Это обозначение приспособлено к пружинной аналогии, используемой в лекции. В стандартных релятивистских обозначениях для струны обычно пишут
\begin{equation}
  T_s=\frac{1}{2\pi\alpha'},
  \label{eq:l05-alpha-prime}
\end{equation}
так что увеличение струнного уровня дает приращение \(m^2\) порядка
\(1/\alpha'\), то есть также порядка натяжения.
\editorialnote{В лекции эффективное \(k\) выбрано так, что \(\sqrt{k}\) играет роль натяжения. Не зависящее от соглашений утверждение состоит в том, что \(T_s\) имеет размерность квадрата массы, а наклон реджевской траектории удовлетворяет соотношению \(\alpha'\propto T_s^{-1}\).}

\subsection{Размерности и интервалы между осцилляторными уровнями}

Положим \(c=\hbar=1\). Тогда энергия имеет размерность обратной длины:
\begin{equation}
  [E]=L^{-1}.
\end{equation}
Поскольку натяжение есть энергия на единицу длины,
\begin{equation}
  [T_s]=L^{-2}=E^2.
  \label{eq:l05-tension-dimension}
\end{equation}
Вот почему струнный спектр естественно разнесен по квадрату массы.

\begin{figure}[H]
  \centering
  \includegraphics[width=0.68\textwidth]{lecture_05_figure_03.png}
  \caption{Естественные единицы на доске: \(c=\hbar=1\), энергия — обратная длина, а частота осциллятора задается масштабом пружины.}
  \lectureframe{00:06:05}
  \label{fig:l05-natural-units-board}
\end{figure}

Обычный осциллятор приводит к тому же выводу с другой стороны. Для пружины с массой \(\mu\),
\begin{equation}
  \omega=\sqrt{\frac{k}{\mu}},
  \qquad
  \Delta E=\hbar\omega.
\end{equation}
Более жесткая пружина имеет более высокую частоту и больший квантовый интервал. Струна содержит бесконечно много таких нормальных мод, но их общий масштаб задается \(T_s\). Как только натяжение известно, известен и интервал между уровнями струнного спектра.

Натяжение имеет также размерность силы, потому что сила, умноженная на расстояние, есть энергия. Представьте, что один конец закреплен, а к другому подвешен груз. Идеальная релятивистская струна с \(E=T_s\ell\) может выдерживать силу, задаваемую \(T_s\), независимо от своей длины. Адронное натяжение, выводимое из наблюдаемых энергий возбуждения, соответствует в этом намеренно причудливом сравнении весу грузовика. К мезону никогда не подвешивали никакого грузовика; измеренный спектр дает натяжение, а натяжение дает механическое сравнение.

Фундаментальные струны были бы гораздо жестче. Если именно они лежат в основе электронов, фотонов и гравитонов, то их первые осцилляторные возбуждения должны находиться намного выше наблюдаемых энергий частиц. В том же причудливом гравитационном эксперименте их натяжение могло бы удерживать нечто галактической массы. Смысл сравнения не в инженерии. Он в том, что огромное натяжение означает и большую энергию возбуждения, и очень малую флуктуацию основного состояния.

\section{Единицы, построенные из универсальных констант}

Какой масштаб достаточно велик? Естественный ориентир — планковский масштаб. Вместо того чтобы просто привести его, давайте его построим.

Метры, килограммы и секунды — превосходные лабораторные единицы, но их происхождение связано с человеческими масштабами. Метр удобен для ткани или веревки; фут назван по части тела; секунда — удобная доля суток. Можно было бы положить радиус протона равным единице и упростить некоторые ядерные формулы, но протон — лишь одна конкретная составная частица. Такой выбор не упростил бы всю физику.

Три размерные константы имеют более универсальный охват:
\begin{itemize}
  \item \(c\) задает предельную скорость для любой материальной системы;
  \item \(\hbar\) входит в соотношения неопределенности любой квантовой
  системы;
  \item \(G\) связывает гравитацию с любой формой массы-энергии.
\end{itemize}
Выбрать единицы, в которых \(c=\hbar=G=1\), — значит построить единицы длины, времени и массы из этих констант.

Их обычные размерности таковы
\begin{equation}
  \begin{aligned}
    [c]&=LT^{-1},
    & [\hbar]&=ML^2T^{-1},\\
    [G]&=L^3M^{-1}T^{-2}.
  \end{aligned}
  \label{eq:l05-dimensions-constants}
\end{equation}
Второе соотношение прямо следует из
\(\Delta x\,\Delta p\gtrsim\hbar\). Последнее следует, например, из
\begin{equation}
  a=\frac{GM}{r^2}.
\end{equation}

\begin{figure}[H]
  \centering
  \begin{minipage}[t]{0.48\textwidth}
    \centering
    \includegraphics[width=\textwidth]{lecture_05_figure_04.png}
    \par\smallskip\footnotesize (a) Размерности \(c\) и \(\hbar\).
  \end{minipage}\hfill
  \begin{minipage}[t]{0.48\textwidth}
    \centering
    \includegraphics[width=\textwidth]{lecture_05_figure_05.png}
    \par\smallskip\footnotesize (b) Размерность \(G\) из ньютоновского ускорения.
  \end{minipage}
  \caption{Элементы анализа размерностей, собранные на занятии.}
  \lectureframe{00:25:43--00:26:40}
  \label{fig:l05-dimensional-analysis-board}
\end{figure}

\subsection{Сначала планковская площадь}

Ищем такие степени \(p,q,r\), при которых
\begin{equation}
  G^p\hbar^q c^r
\end{equation}
имеет размерность \(L^2\). Выбор площади вначале позволяет избежать раннего извлечения квадратного корня и одновременно предвосхищает ту особую роль, которую площадь позже сыграет в гравитации. Из \eqref{eq:l05-dimensions-constants}, сокращая массу и время, получаем
\begin{equation}
  -p+q=0,
  \qquad
  -2p-q-r=0.
\end{equation}
Для оставшейся степени длины должно выполняться
\begin{equation}
  3p+2q+r=2.
\end{equation}
Следовательно, \(q=p\), \(r=-3p\) и \(p=1\). Планковская площадь равна
\begin{equation}
  \ell_P^2=\frac{G\hbar}{c^3},
  \label{eq:l05-planck-area}
\end{equation}
и, следовательно,
\begin{equation}
  \ell_P=\sqrt{\frac{G\hbar}{c^3}},
  \qquad
  t_P=\frac{\ell_P}{c},
  \qquad
  m_P=\sqrt{\frac{\hbar c}{G}}.
  \label{eq:l05-planck-units}
\end{equation}

\begin{figure}[H]
  \centering
  \includegraphics[width=0.78\textwidth]{lecture_05_figure_06.png}
  \caption{Завершенный размерностный вывод для планковской площади и длины.}
  \lectureframe{00:31:34}
  \label{fig:l05-planck-area-board}
\end{figure}

С точностью до порядка величины,
\begin{equation}
  \begin{aligned}
    \ell_P&\sim10^{-35}\,\mathrm{m},
    &t_P&\sim10^{-43}\,\mathrm{s},\\
    m_P&\sim10^{-8}\,\mathrm{kg},
    &m_Pc^2&\sim10^{19}\,\mathrm{GeV}.
  \end{aligned}
  \label{eq:l05-planck-numbers}
\end{equation}

\begin{figure}[H]
  \centering
  \includegraphics[width=0.72\textwidth]{lecture_05_figure_07.png}
  \caption{Три планковских масштаба, записанные с точностью до порядка величины.}
  \lectureframe{00:35:12}
  \label{fig:l05-planck-values-board}
\end{figure}

Длина и время невообразимо малы, тогда как масса по меркам физики частиц удивительно макроскопична, примерно как масса мелкой пылинки. Ее энергия покоя имеет порядок нескольких гигажоулей, сопоставимый с химической энергией в баке бензина. Наблюдаемая Вселенная, напротив, имеет размер примерно \(10^{60}\) планковских длин и существует примерно \(10^{60}\) планковских времен. Эти сравнения призваны придать единицам физическую весомость, а не улучшить их численную точность.

Если характерная длина струны \(\ell_s=\sqrt{\alpha'}\) близка к
\(\ell_P\), то квантовый размер струны близок к \(\ell_P\), период ее колебаний — к \(t_P\), а ее первое массивное возбуждение — к
\(m_P\). В лекции допускаются множители 10, 100 и больше; это отождествление — ожидание, а не размерностное тождество.

\begin{classroomqa}
\audiencequestion{Является ли масса планковского масштаба приращением относительно основного состояния частицы? \lecturetimestamp{00:38:37}}
\lecturerresponse{Равномерно разнесенной величиной является квадрат массы. Если основное состояние легкое, то у первого струнного возбуждения \(m^2\) больше на величину порядка \(m_P^2\), так что его масса и требуемая энергия в системе центра масс имеют порядок \(m_P\).}
\end{classroomqa}

\section{Насколько близки должны быть эти два масштаба?}

\begin{classroomqa}
\audiencequestion{Анализ размерностей дает \(\ell_P\), но почему эта величина должна быть длиной струны? \lecturetimestamp{00:40:21}}
\lecturerresponse{Это физическое предположение, подкрепляемое гравитацией и отсутствием наблюдаемой струнной структуры, а не следствие одного лишь анализа размерностей. Длина струны может быть больше планковской длины, а это означает, что ее энергетический масштаб ниже. В лекции утверждается лишь, что она не должна быть обычной длиной в физике частиц и что если опустить ее намного ниже планковской длины, то это приведет к гравитационному коллапсу.}
\end{classroomqa}

Связь между двумя масштабами в действительности зависит от константы связи. В полной компактификации \(\ell_P\), \(\ell_s\), константа струнной связи и объем компактификации входят совместно. Универсального уравнения \(\ell_s=\ell_P\) не существует. Из лекции сохраняется аргумент о масштабах: струнные эффекты должны проявляться лишь на очень высоких энергиях, и субпланковскую локализацию нельзя обсуждать, пренебрегая гравитацией.
\editorialnote{В современных струнных конструкциях не обязательно точно отождествлять струнную длину с четырехмерной планковской длиной. Их соотношение зависит от константы струнной связи и данных компактификации.}

\begin{classroomqa}
\audiencequestion{Если считать планковскую длину наименьшей осмысленной длиной, является ли планковская масса также наименьшей возможной массой? \lecturetimestamp{00:43:27}}
\lecturerresponse{Нет. Обычные частицы гораздо легче. Планковская масса — это масштаб, на котором компактный объект может иметь гравитационный радиус, сравнимый с его квантовой длиной волны. На эвристическом языке лекции это масштаб самой легкой полуклассической черной дыры со временем жизни порядка планковского времени; последнее утверждение не следует принимать за точное описание полностью квантового объекта.}
\end{classroomqa}

\subsection{Что могут сказать нам бегущие константы связи}

Имеются косвенные свидетельства того, что обычная квантовая теория поля остается справедливой на масштабах, далеко выходящих за пределы прямой досягаемости ускорителей. Три калибровочные константы связи Стандартной модели измеряются при доступных энергиях, и их зависимость от масштаба можно экстраполировать. На схематическом графике из лекции константы связи \(U(1)\), \(SU(2)\) и \(SU(3)\) сближаются друг с другом на высоком масштабе объединения, возможно на несколько порядков ниже планковской энергии.

\begin{figure}[H]
  \centering
  \includegraphics[width=0.70\textwidth]{lecture_05_figure_08.png}
  \caption{Схематический график бегущих констант связи. Три калибровочные кривые сближаются друг с другом до того, как возрастающая гравитационная сила становится величиной порядка единицы.}
  \lectureframe{00:47:18}
  \label{fig:l05-running-couplings-board}
\end{figure}

Это экстраполяция, а не прямой эксперимент в точке пересечения. Дополнительные частицы или пороги между лабораторными энергиями и энергиями объединения могут изменить наклоны. Осторожный вывод состоит лишь в том, что известная теория поля не обнаруживает признаков подструктуры струнного размера на обычных расстояниях и что правдоподобный струнный масштаб лежит где-то в огромном интервале между выводимым масштабом объединения и планковским масштабом.

\begin{classroomqa}
\audiencequestion{Какие именно три константы связи экстраполируются? \lecturetimestamp{00:48:28}}
\lecturerresponse{Это калибровочные константы связи, соответствующие \(U(1)\), \(SU(2)\) и \(SU(3)\): гиперзарядному, слабому и сильному взаимодействиям. Электромагнетизм возникает из электрослабого смешивания, а не является в этом учете независимым четвертым калибровочным фактором.}
\end{classroomqa}

\begin{classroomqa}
\audiencequestion{Растет ли сама гравитационная постоянная \(G\) по мере того, как мы зондируем более высокие энергии? \lecturetimestamp{00:49:30}}
\lecturerresponse{\(G\) имеет размерность, поэтому сама по себе это не та величина, которую следует сравнивать с безразмерной калибровочной константой связи. При характерной энергии \(M\) безразмерная гравитационная сила схематически ведет себя как \(GM^2/(\hbar c)\). Эта комбинация растет с энергией и на планковском масштабе становится величиной порядка единицы.}
\end{classroomqa}

Это сравнение становится наглядным из законов сил:
\begin{equation}
  F_{\mathrm{em}}\sim\frac{e^2}{r^2},
  \qquad
  F_{\mathrm{grav}}\sim\frac{GM^2}{r^2}.
  \label{eq:l05-force-laws}
\end{equation}
Таким образом, \(e^2\) сравнивают не с \(G\), а с \(GM^2\) в естественных единицах.

\begin{figure}[H]
  \centering
  \includegraphics[width=0.74\textwidth]{lecture_05_figure_09.png}
  \caption{Законы электромагнитной и гравитационной сил, используемые для определения безразмерной гравитационной силы.}
  \lectureframe{00:50:50}
  \label{fig:l05-force-laws-board}
\end{figure}

\section{Почему прямая струнная спектроскопия недостижима}

Как установить экспериментально, что электрон является струной? Мы сделали бы то же, что делала адронная спектроскопия: ударили бы по нему достаточно сильно, чтобы возбудить его внутренние моды, и искали бы башню состояний, массы которых следуют струнной схеме. Если первый массивный уровень лежит около \(10^{19}\,\mathrm{GeV}\), это далеко за пределами возможностей любого прямого ускорителя.

Полная энергия — не единственное препятствие. Ее нужно сосредоточить при чрезвычайно малом параметре удара. Здесь вступают в дело черные дыры.

\begin{classroomqa}
\audiencequestion{На этих энергиях начнет ли ускоритель производить малые черные дыры? \lecturetimestamp{00:50:57}}
\lecturerresponse{Таково эвристическое ожидание, как только энергия в системе центра масс сосредоточивается внутри своего гравитационного радиуса. Звездная черная дыра образуется, когда гравитационная потенциальная энергия сжимает массивную звезду. Микроскопическая же образовалась бы из кинетической энергии, сфокусированной в крошечной области столкновения. В четырехмерной гравитации требуемый масштаб — планковский.}
\end{classroomqa}

Для ориентира прирост энергии в линейном ускорителе примерно пропорционален его длине. В лекции в качестве простого ориентира масштаба SLAC используется оценка: около двух миль на примерно \(100\,\mathrm{GeV}\). Достижение \(10^{19}\,\mathrm{GeV}\) при том же градиенте увеличило бы длину примерно в \(10^{17}\) раз, дав машину галактического масштаба.

Даже такой машине потребовалась бы светимость. Пучки не наводят попарно с точностью до планковской длины; пересекаются огромные числа частиц, так что немногие сталкиваются при нужном параметре удара. Умножение планковской энергии на огромную частоту столкновений дает столь же абсурдное требование к мощности. Оценки намеренно грубы, но их вывод устойчив: прямая струнная спектроскопия на планковском масштабе — не тот эксперимент, который можно сколько-нибудь реалистично спроектировать.

Косвенные свидетельства должны были бы носить более структурный характер: теория, которая выводит наблюдаемые низкоэнергетические частицы и взаимодействия с достаточной определенностью, чтобы ее можно было проверить, или эффекты компактификации, сохраняющиеся далеко ниже струнного масштаба. В лекции не утверждается, что такая решающая конструкция уже имеется.

\subsection{Дополнительные измерения и предел локализации}

\begin{classroomqa}
\audiencequestion{Должны ли дополнительные измерения тоже иметь размеры, близкие к планковской длине? \lecturetimestamp{00:56:23}}
\lecturerresponse{В этой элементарной картине ожидается, что они будут чрезвычайно малы, хотя аргумент не фиксирует их точно. Если бы компактное измерение было достаточно большим, чтобы его можно было исследовать при меньшей энергии, новые Kaluza--Klein-моды и более высокоразмерное поведение уже изменили бы наблюдаемую нами теорию поля. Поэтому ее успешная экстраполяция накладывает верхнюю границу на размеры таких измерений.}
\end{classroomqa}

\begin{classroomqa}
\audiencequestion{Почему струна или компактное измерение не могли бы быть на многие порядки меньше, чем \(\ell_P\)? \lecturetimestamp{00:57:48}}
\lecturerresponse{Разрешение меньшего расстояния требует большего импульса. Выше планковского масштаба энергия, необходимая для локализации, имеет гравитационный радиус, сравнимый с исследуемой областью, а затем и превышающий ее. Эксперимент образует черную дыру прежде, чем разрешит меньшую обычную струну.}
\end{classroomqa}

Принцип неопределенности дает оценку. Локализация до \(\Delta x\) требует

\begin{classroomqa}
\audiencequestion{После того как выбраны планковская длина и планковское время, как физически определяется планковская масса? \lecturetimestamp{00:59:04}}
\lecturerresponse{Можно использовать либо \(\hbar\), либо \(G\), чтобы перевести масштаб длины в энергию. Принцип неопределенности говорит, что локализация двух частиц внутри \(\ell_P\) требует импульса порядка \(\hbar/\ell_P\), а значит, энергии порядка \(m_Pc^2\). Эквивалентно, это масса, при которой квантовая и гравитационная длины становятся сравнимыми.}
\end{classroomqa}

\begin{equation}
  \Delta p\gtrsim\frac{\hbar}{\Delta x},
  \qquad
  E\sim c\,\Delta p.
  \label{eq:l05-localization}
\end{equation}
При \(\Delta x=\ell_P\) эта энергия равна
\begin{equation}
  E\sim\frac{\hbar c}{\ell_P}=m_Pc^2.
\end{equation}
В то же время радиус Шварцшильда
\begin{equation}
  r_s\sim\frac{GE}{c^4}
\end{equation}
имеет порядок \(\ell_P\) с точностью до численных констант. Квантовая локализация и гравитационный коллапс сходятся на одном и том же масштабе.

\begin{figure}[H]
  \centering
  \includegraphics[width=0.74\textwidth]{lecture_05_figure_10.png}
  \caption{Формулы для планковской длины и массы, к которым вновь обращаются в аргументе о локализации.}
  \lectureframe{01:00:25}
  \label{fig:l05-localization-board}
\end{figure}

На этом аргумент о выборе масштаба завершен.

\section{Основное состояние должно лежать ниже фотона}

После перерыва возможны два пути: начать конформные преобразования или изложить более старый и более простой путь к критической размерности. Мы выбираем второй, но с предупреждением. Простой аргумент не так уж прост и не является самой глубокой причиной. Тем не менее он верен и близок к историческому открытию.

Вернемся к открытой бозонной струне. Однократное действие поперечного оператора рождения создает два поляризационных состояния. Поскольку продольного третьего состояния нет, такое возбуждение должно быть безмассовым вектором. В единицах, где один квант моды-\(1\) увеличивает \(m^2\) на единицу,
\begin{equation}
  m_{\mathrm{ground}}^2=-1,
  \qquad
  m_{\mathrm{first}}^2=-1+1=0.
  \label{eq:l05-open-intercept}
\end{equation}
Тахионное основное состояние остается проблемой бозонной теории; это признается и откладывается в сторону. Непосредственный вопрос состоит в том, откуда берется отрицательная единица в \eqref{eq:l05-open-intercept}.

Каждая мода струны является гармоническим осциллятором. При \(\hbar=1\) ее энергия вакуума равна
\begin{equation}
  E_{0,n}=\frac{\omega_n}{2}.
\end{equation}
Для струны \(\omega_n=n\), так что одна поперечная координата формально дает вклад
\begin{equation}
  E_0^{(1)}
  =\frac12\sum_{n=1}^{\infty}n.
  \label{eq:l05-formal-zero-point}
\end{equation}

\begin{figure}[H]
  \centering
  \includegraphics[width=0.70\textwidth]{lecture_05_figure_11.png}
  \caption{Сумма нулевой энергии \(\sum n/2\), записанная после того, как были напомнены частоты мод.}
  \lectureframe{01:06:18}
  \label{fig:l05-zero-point-board}
\end{figure}

Теперь трудность становится особенно острой. Каждый член в
\eqref{eq:l05-formal-zero-point} положителен, и сумма расходится. Как она может давать конечный отрицательный интерсепт?

\section{Прежде чем читать спектр, вычтем движение}

Прежде чем регуляризовать сумму, напомним, почему внутренняя энергия отождествлялась с квадратом массы. Для частицы, несущей большой продольный импульс
\(P>0\),
\begin{equation}
  E=\sqrt{P^2+m^2}
   =P+\frac{m^2}{2P}
   +O(P^{-3}).
  \label{eq:l05-large-p-energy}
\end{equation}

\begin{figure}[H]
  \centering
  \includegraphics[width=0.70\textwidth]{lecture_05_figure_12.png}
  \caption{Разложение при больших
  \(P\), отделяющее вклад общего импульса от внутреннего массового члена.}
  \lectureframe{01:07:35}
  \label{fig:l05-fast-energy-board}
\end{figure}

Первый член — это неинтересное общее движение вдоль оси \(z\). Он сохраняется и одинаков для всех внутренних состояний с одним и тем же \(P\). Вычитание этого вклада и последующее перемасштабирование устраняют кинематическое замедление времени:
\begin{equation}
  \begin{aligned}
    E-P&=\frac{m^2}{2P}+O(P^{-3}),\\
    2P(E-P)&=m^2+O(P^{-2}).
  \end{aligned}
  \label{eq:l05-internal-energy}
\end{equation}

Атом дает наглядную физическую картину. Подвергнем атом водорода лоренцеву бусту. Его полная энергия возрастает, но разность энергий двух внутренних уровней, измеряемая в лабораторном времени, уменьшается, потому что ход часов атома замедлен. Множитель \(1/P\) в \eqref{eq:l05-large-p-energy} и есть это подавление. Умножение на \(2P\) восстанавливает инвариантную разность квадратов масс.

Поэтому необработанная энергия основного состояния не обязана буквально быть равной \(-1\). Она может содержать большой не зависящий от состояния вклад, связанный с общей продольной энергией, плюс конечный внутренний остаток. В лекции соотношение после буста переписывается так, чтобы явно показать обе части, прежде чем вернуться к сумме осцилляторов.

\begin{figure}[H]
  \centering
  \includegraphics[width=0.76\textwidth]{lecture_05_figure_13.png}
  \caption{Инвариантный массовый член, выписанный рядом с большим
  не зависящим от состояния продольным вкладом.}
  \lectureframe{01:11:34}
  \label{fig:l05-infinite-piece-board}
\end{figure}

Бесконечные энергии вакуума хорошо знакомы по квантовой теории поля. Часто отделяют расходящуюся не зависящую от состояния константу от наблюдаемых разностей энергий. В лекции это сопоставляется с энергией нулевых колебаний электромагнитного поля в комнате, и сама процедура прямо называется неудовлетворительной, хотя и стандартной. Здесь мы следуем той же эвристике: выделяем расходящуюся часть и спрашиваем, сохраняется ли определенный конечный остаток.

\section{Регуляризация суммы}

Введем положительный регулятор \(\epsilon\):
\begin{equation}
  S_1(\epsilon)
  =\sum_{n=1}^{\infty}n e^{-n\epsilon}.
  \label{eq:l05-regulated-sum}
\end{equation}
При любом \(\epsilon>0\) экспонента в конце концов пересиливает множитель
\(n\), так что сумма сходится. Предел \(\epsilon\) к нулю снимает
отсечку.

\begin{figure}[H]
  \centering
  \includegraphics[width=0.78\textwidth]{lecture_05_figure_14.png}
  \caption{Экспоненциальный регулятор, примененный почленно к
  \(1+2+3+\cdots\).}
  \lectureframe{01:16:42}
  \label{fig:l05-regulator-board}
\end{figure}

Множитель \(n\) можно получить дифференцированием:
\begin{equation}
  S_1(\epsilon)
  =
  -\frac{d}{d\epsilon}
  \sum_{n=1}^{\infty}e^{-n\epsilon}.
  \label{eq:l05-derivative-trick}
\end{equation}

\begin{figure}[H]
  \centering
  \includegraphics[width=0.76\textwidth]{lecture_05_figure_15.png}
  \caption{Дифференцирование по параметру отсечки дает множитель
  \(-n\).}
  \lectureframe{01:18:05}
  \label{fig:l05-derivative-board}
\end{figure}

Оставшийся ряд — геометрический:
\begin{equation}
  S_0(\epsilon)
  \equiv
  \sum_{n=1}^{\infty}e^{-n\epsilon}
  =
  \frac{e^{-\epsilon}}{1-e^{-\epsilon}}
  =
  \frac{1}{e^\epsilon-1}.
  \label{eq:l05-geometric-sum}
\end{equation}

\begin{figure}[H]
  \centering
  \includegraphics[width=0.68\textwidth]{lecture_05_figure_16.png}
  \caption{Регуляризованный геометрический ряд в замкнутом виде.}
  \lectureframe{01:19:12}
  \label{fig:l05-geometric-board}
\end{figure}

Разложим при малом \(\epsilon\):
\begin{equation}
  S_0(\epsilon)
  =
  \frac{1}{\epsilon}
  -\frac12
  +\frac{\epsilon}{12}
  -\frac{\epsilon^3}{720}
  +O(\epsilon^5).
  \label{eq:l05-bernoulli-expansion}
\end{equation}
На доске это получается раздельным разложением числителя и знаменателя, вынесением одного множителя \(\epsilon\) и последующим разложением обратной величины. В этом вычислении по ходу лекции несколько раз проверяются знаки; компактное выражение \eqref{eq:l05-bernoulli-expansion} — его окончательный аккуратный результат.

\begin{figure}[H]
  \centering
  \includegraphics[width=0.78\textwidth]{lecture_05_figure_17.png}
  \caption{Разложения числителя и знаменателя до упрощения обратной
  величины.}
  \lectureframe{01:23:10}
  \label{fig:l05-expansion-board}
\end{figure}

Применяя \(-d/d\epsilon\), получаем
\begin{equation}
  S_1(\epsilon)
  =
  \frac{1}{\epsilon^2}
  -\frac{1}{12}
  +\frac{\epsilon^2}{240}
  +O(\epsilon^4).
  \label{eq:l05-regulated-result}
\end{equation}

\begin{figure}[H]
  \centering
  \includegraphics[width=0.78\textwidth]{lecture_05_figure_18.png}
  \caption{Расходящийся член \(1/\epsilon^2\) и конечный
  остаток \(-1/12\) в конце вычисления с регулятором.}
  \lectureframe{01:29:48}
  \label{fig:l05-minus-twelve-board}
\end{figure}

Часто этот результат в шутку резюмируют как
\(1+2+3+\cdots=-1/12\). Это не обычная сходимость. Уравнение
\eqref{eq:l05-regulated-result} говорит, что именно этот регулятор дает
расходящийся член плюс не зависящую от регулятора конечную часть:
\begin{equation}
  \underbrace{\frac{1}{\epsilon^2}}_{\text{divergent}}
  \;+\;
  \underbrace{\left(-\frac{1}{12}\right)}_{\text{finite}}
  \;+\;O(\epsilon^2).
\end{equation}
В лекции первый член поглощается в уже имеющийся
не зависящий от состояния вклад, а второй оставляется. Также приводится
предупреждение, приписываемое Паули: то, что величина
бесконечна, само по себе не делает ее нулем. Согласованная теория должна
объяснять, почему такое вычитание допустимо.

Если восстановить множитель одной второй в
\eqref{eq:l05-formal-zero-point}, то одна поперечная координата дает вклад
\begin{equation}
  E_0^{(1)}=-\frac{1}{24}.
  \label{eq:l05-one-transverse-zero-point}
\end{equation}

\begin{figure}[H]
  \centering
  \includegraphics[width=0.68\textwidth]{lecture_05_figure_19.png}
  \caption{Множитель одна вторая превращает конечное
  \(-1/12\) в \(-1/24\) на каждую поперечную координату.}
  \lectureframe{01:33:42}
  \label{fig:l05-minus-twenty-four-board}
\end{figure}

\section{Подсчет поперечных направлений}

Пусть \(D_\perp\) — число поперечных координат струны. Их полный
вакуумный вклад равен
\begin{equation}
  E_0^{\mathrm{open}}
  =-\frac{D_\perp}{24}.
  \label{eq:l05-open-vacuum}
\end{equation}
Чтобы первый осциллятор открытой струны поднял основное состояние с
\(-1\) до безмассового уровня, необходимо
\begin{equation}
  -\frac{D_\perp}{24}=-1.
\end{equation}
Следовательно
\begin{equation}
  D_\perp=24.
\end{equation}
Это направления, перпендикулярные большому продольному импульсу.
Прибавляя это продольное пространственное направление и время, получаем
\begin{equation}
  D=D_\perp+2=26.
  \label{eq:l05-critical-dimension}
\end{equation}

\begin{classroomqa}
\audiencequestion{Когда этот аргумент вошел в теорию
струн? \lecturetimestamp{01:35:27}}
\lecturerresponse{Он относится к первоначальной эпохе адронных
струн, примерно к 1969 году и началу 1970-х. Траектории Редже уже были
известны экспериментально. Когда осцилляторную модель стали исследовать
саму по себе, оказалось, что первое возбужденное состояние имеет две, а не три векторные
поляризации, и потому должно быть безмассовым. Требуемый интерсепт затем
быстро привел к двадцати четырем поперечным осцилляторам.}
\end{classroomqa}

Тогда это казалось почти слишком странным, чтобы в это поверить. Почему регуляризованная вакуумная сумма должна знать размерность пространства-времени? И все же все сошлось: безмассовый вектор, интерсепт и число поперечных осцилляторов указывали на одно и то же число.

\subsection{Замкнутая струна проверяет подсчет}

Замкнутой струне нужны одна леводвижущаяся и одна праводвижущаяся мода, чтобы достичь своего первого физического безмассового уровня. Поэтому ее основное состояние должно начинаться на две единицы ниже нуля.

\begin{classroomqa}
\audiencequestion{Если двадцать четыре поперечных направления дают только \(-1\),
должна ли замкнутая струна тогда существовать в сорока восьми поперечных
направлениях? \lecturetimestamp{01:37:47}}
\lecturerresponse{Нет. Одно и то же пространство-время содержит и открытые, и замкнутые
струны. У замкнутой струны уже вдвое больше осцилляторных мод: по одной башне
мод в каждом направлении мирового листа. Каждая башня дает вклад
\(-24/24=-1\), так что вместе они дают \(-2\).}
\end{classroomqa}

Итак,
\begin{equation}
  E_0^{\mathrm{closed}}
  =
  2\left(-\frac{D_\perp}{24}\right)
  =-2
  \qquad (D_\perp=24).
\end{equation}
Добавление одного осциллятора из каждого сектора поднимает состояние на две единицы и приводит к нулевому квадрату массы, точно туда, где в предыдущей лекции появился гравитонный мультиплет.

Это совпадение делает аргумент более правдоподобным, но глубоким объяснением оно все еще не является. В современной квантизации константа нормального упорядочения, замыкание алгебры Лоренца или Вирасоро и устранение конформной аномалии входят в одну согласованную структуру. Эта структура фиксирует и интерсепт, и критическую размерность; конечный остаток после отсечки — исторический путь к ней.
\editorialnote{Вычисление с регулятором — эвристический взгляд на константу нормального упорядочения. Более глубокий вывод \(D=26\) требует конформной инвариантности мирового листа, или, что эквивалентно, устранения конформной аномалии и согласованной алгебры связей.}

Теперь мы перейдем к более глубокому подходу. Чтобы начать, нужна лишь небольшая часть комплексного анализа: комплексные координаты, аналитические функции и уравнения Коши--Римана.
Обещанный физический образ — это Turkish taffy---ее тянут в одну сторону, растягивают в другую и придают ей новую форму, не разрушая существенной структуры. Стоящая за этим образом математика — конформная инвариантность.